\section{Limitations and Future Work}
%This paper has two limitations. Firstly, FunQA's input is limited to raw videos, which restricts the model's ability to process and analyze the content effectively. Secondly, our experimental results reveal a relatively low overall score. Therefore, future work could involve performing additional annotation of the video, such as panoramic segmentation of images, to create a more extensive and comprehensive VideoQA dataset. Additionally, exploring the development of novel evaluation metrics is a promising research direction in the VideoQA field. Secondly,  exploring the development of novel evaluation metrics holds promise as a research avenue within the VideoQA field.\\
%In conclusion, as a new benchmark, FunQA paves the way for exploring interesting videos and delving into deeper human understanding. By tackling the aforementioned limitations, future research can further enhance the performance of models in video comprehension tasks.
This paper has two limitations. \textbf{1)}\quad Current FunQA dataset mainly includes video-level data and annotations, but denser annotations (e.g., like PVSG~\cite{yang2023panoptic}) can be developed to explore more possibilities of video reasoning. Examples include detailed spatial and temporal annotations, such as captions corresponding with specific time axes and annotations of object level. \textbf{2)}\quad The raw annotations are completed by the annotator in Chinese. In the process of translating into English, we first use GPT to polish and supplement Chinese annotations, making the text as thorough as possible, as short or incomplete text may result in misunderstanding. However, there may still be differences due to cultural diversities between the two languages.

In the future, we expect to expand our dataset with denser and more diverse annotations. Also, new metrics will be explored to better evaluate models' performance, especially in open-ended questions which lack in-depth metrics. Finally, we hope to provide directions for advancing models toward deeper video reasoning. Specific ideas include that Otter can sample more frames from the input video to improve causal reasoning capacity. Furthermore, in the future, when refining tuning datasets such as MIMIC-IT~\cite{li2023mimic} for visual language models, it may be beneficial to incorporate the FunQA training set. However, to ensure fair comparisons and prevent data leakage, it is advised that future research does not utilize the FunQA testing set.